\documentclass[12pt]{article}

\usepackage{sbc-template}
\usepackage{graphicx,url}
\usepackage[utf8]{inputenc}
%\usepackage[english]{babel}
\usepackage{cite}
\usepackage{overpic}
\usepackage{amsmath, amssymb, float}
\usepackage[inline]{enumitem}
\usepackage{xcolor}
%\usepackage{hyperref}

\sloppy

\renewcommand{\thefootnote}{\fnsymbol{footnote}}

\newcommand{\secrefplain}[1]{Section~\ref{#1}}
\newcommand{\myvector}[1]{{\boldsymbol{\mathbf{ #1 }}}}

\hyphenpenalty=10000
\tolerance=2000

\title{Efficient Computation of the Directional Extremal Boundary\\
of a Union of Equal-Radius Circles}

\author{Caio Grimm}

\address{Instituto de Computação - UFAM\\
\email{caiogrimm@gmail.com}}

\begin{document}
	
\maketitle

\begin{abstract}
This paper focuses on computing the directional extremal boundary
of a union of equal-radius circles. We introduce an efficient algorithm
that accurately determines this boundary by analyzing the intersections and
dominant relationships among the circles. The algorithm has time complexity
of $O(n \log n)$.
\end{abstract}

\section{Introduction}
Computing the directional extremal boundary of a union of equal-radius circles involves
identifying the outermost points in a specified direction that form the outer edge of the
combined circle regions. This task is essential in a wide range of applications where each
circle may represent a measurement, area of influence, or coverage zone.

Without loss of generality, this paper focuses on computing the upper boundary of
the union of unit circles, as this can be readily extended to circles of any equal radius by
applying uniform scaling and rotation. In this paper, we present an efficient algorithm for
determining the upper boundary of overlapping unit circles. The method begins by sorting
all circle centers according to their horizontal positions. Once sorted, the circles are
processed from left to right to build the boundary incrementally, accounting for overlaps
and overshadowing. Specifically, it identifies the exact horizontal positions where one
circle's boundary overtakes or merges into another's. By maintaining a dynamic list of
circles that currently define the boundary, the algorithm eliminates any circles that no
longer contribute to the highest points, thereby preserving only those that shape the overall
upper boundary.

One practical application of this algorithm is the precise localization of extremal
points along elongated structures such as rail tracks or pipelines. In such cases, each
circle may represent a measurement along the structure, and computing the circles'
directional extremal boundary enables accurate identification of the outermost positions
in the specified direction.

Another application arises in robotic workspace modeling, where a robotic arm
moves along rails in one direction and extends in another. Here, circles can represent
obstacles and safety zones: the circle centers reflect object positions, while radii model
object sizes, arm dimensions, and required safety distances. The algorithm can then
determine safe regions in which the arm may operate without collisions.

\section{Problem Statement}

\section{Algorithm Description}

\section{Proof of Algorithm Correctness}

\section{Complexity Analysis}

\section{Binary Search for Finding the Upper Value}

\section{Conclusion}

\section*{Acknowledgments}

\section*{Appendix: Horizontal Transition Position of Two Unit Circles}

\bibliographystyle{sbc}
\bibliography{sbc-template}

\end{document}
